\section{Metodología de Gestión de Proyectos}
La dirección utiliza PMBOK como marco base, conforme al artículo 4.3 y RT-19.01. Synaptix adopta gestión híbrida: líneas base y control de compromisos, con revisión iterativa de productos. El procedimiento de adaptación publicado por PMI en 2021 para la séptima edición propone elegir el enfoque, ajustarlo a la organización y al proyecto, y mejorar continuamente. Las reglas operativas siguientes son decisiones de esta oferta, sujetas a las bases. \parencites[art.~4.3, pp.~5--6]{sd6-ba}[RT-19.01, p.~33]{sd6-bt}[pp.~1--7]{sd6-pmi-adaptacion}

\subsection{Enfoque aplicado y líneas base}
El enfoque híbrido combina hitos invariables e instalación física con procesos que requieren validación frecuente de las personas usuarias. PMI reconoce la combinación de prácticas predictivas y ágiles; Synaptix la concreta mediante líneas base y un flujo de construcción que permite precisar requisitos sin ampliar unilateralmente el contrato. \parencite{sd6-pmi-hibrido}

El Jefe de Proyecto constituirá la oficina de gestión y consolidará alcance, requisitos, cronograma, recursos, calidad, riesgos, comunicaciones, interesados y adquisiciones. Conservará códigos de SD7 y la relación requisito--entregable--paquete--prueba. Enunciado, EDT y diccionario definen alcance; red y calendario definen programación. La aprobación de las líneas base se documentará con el CLIENTE: la formulación de la oferta no acredita esa aprobación. La estimación ascendente, el contraste UCP y CPM/PERT de SD7 son sus instrumentos. Los puntos funcionales no se vuelven a sumar a las horas distribuidas, y una duración PERT no equivale a probabilidad de entrega sin sustento. \parencites[RT-19.03--05, p.~33]{sd6-bt}[art.~72, p.~38]{sd6-ba}

Los meses 13--15 combinan marcha blanca E1 y construcción E2; los meses 19--20, marcha blanca E2 y preparación de Operación. La selección quincenal comprueba la capacidad por persona en todos sus frentes, incluidos comités y soporte. Un recurso compartido no cuenta como dos personas. Un incremento construido no habilita por sí mismo la producción.

\subsection{Gobierno, interesados y comunicaciones}
El Jefe de Proyecto tendrá dedicación exclusiva en implementación y facultades de ejecución. La Contraparte Técnica validará productos; el Administrador del Contrato ejercerá sus atribuciones contractuales. Arquitectura, Seguridad, Datos, Calidad, Integración y Operación deciden en sus competencias. El CLIENTE conserva la aceptación. La nómina y acreditaciones corresponden a SD12/T-8; la capacidad se programa en SD7/T-15. \parencites[arts.~70--71, p.~37]{sd6-ba}[cap.~19, pp.~33--34]{sd6-bt}

La Tabla~\ref{tab:sd6-gobierno} resume el gobierno obligatorio. La oficina prepara expedientes con decisión requerida, alternativas, impacto y responsable. Las actas se emiten dentro de dos días hábiles, con acuerdos, responsables y plazos. Un comité técnico no sustituye la aprobación contractual de alcance.
\begin{table}[htbp]
\caption{Cadencias de gobierno del proyecto}\label{tab:sd6-gobierno}
\begin{tabularx}{\linewidth}{L{.23\linewidth}L{.22\linewidth}Y}
\hline
\EncabezadoTabla{Instancia} & \EncabezadoTabla{Frecuencia} & \EncabezadoTabla{Resultado principal}\\\hline
\rowcolor{SynSurface} Comité Ejecutivo & Mensual & Decisiones estratégicas, cambios de alcance y riesgos mayores.\\\hline
Comité de Proyecto & Quincenal & Entregables, desviaciones y decisiones operativas.\\\hline
\rowcolor{SynSurface} Comité de Arquitectura & Mensual & Arquitectura, seguridad y deuda técnica.\\\hline
Comité de Operación & Mensual desde mes 13 & Servicio, incidentes, problemas y mejora.\\\hline
\rowcolor{SynSurface} Seguimiento & Semanal & Coordinación, impedimentos y compromisos.\\\hline
\end{tabularx}
\FuenteTabla{Bases Administrativas, art.~71, pp.~37--38; Bases Técnicas Transversales, RT-19.09, p.~34.}
\end{table}
Estas instancias separan autorización contractual y revisión de ingeniería. Los participantes y productos están en el Formulario T-9, \ArchivoTNueve. SRE participa desde el mes 6; el Comité de Operación comienza en 13 sin adelantar la Operación contractual de 21--56.

T-9 concreta calendario y convocatoria, secretaría y representaciones suplentes acreditadas, revisión del borrador en el primer día hábil y publicación del acta dentro de dos días hábiles de cada sesión. Una discrepancia se registra sin atribuir aprobación por silencio ni diferir la publicación; las correcciones mantienen versiones y comunicación a participantes. Proyecto sigue cada acuerdo hasta su evidencia de cierre. La preparación, asistencia y revisión se imputan una sola vez a los paquetes y dedicaciones de SD1/SD7/T-15, con balance por persona, cobertura de ausencias y relevo de mesa; no se usan como capacidad adicional las dos funciones de un mismo titular. El calendario comprobará los solapamientos 13--15 y 19--20 y mantendrá el Comité de Operación y la cobertura inicial desde el mes 13, y la fase contractual de Operación de ambos alcances desde el primer día del mes 21 hasta el 56. \parencites[art.~71, pp.~37--38]{sd6-ba}[RT-19.09, p.~34]{sd6-bt}

El registro de interesados identifica proceso, autoridad, impacto, disponibilidad y canal. La Contraparte coordina sesiones breves con administración, torre, portería, varadero y escuela; Funcional consolida observaciones y conflictos. Las validaciones se agrupan por proceso para que las 39 personas de invierno no deban actuar como un equipo TI permanente. Salud, menores y expedientes personales se muestran sólo a perfiles autorizados. Documentación, actas, riesgos y cambios permanecen en el espacio colaborativo accesible al CLIENTE. \parencites[numerales 13.1--13.3, pp.~28--30]{sd6-c07}[RT-19.05/08, pp.~33--34]{sd6-bt}

\subsection{Control integrado de cambios y adquisiciones}
Una observación se clasifica como defecto, precisión dentro del alcance o modificación. Un defecto se corrige; una precisión conserva obligación y recepción. Si cambia alcance, plazo, costo, riesgo o servicio se tramita solicitud contractual. La Figura~\ref{fig:sd6-cambios} muestra su decisión y registros. \parencite[art.~72, p.~38]{sd6-ba}
\begin{figure}[htbp]
\centering
\begin{tikzpicture}[
  caja/.style={draw=SynLine,fill=SynSurface,rounded corners=1mm,align=left,text width=12.4cm,inner sep=3mm,font=\fontsize{9}{12}\selectfont},
  enlace/.style={-{Stealth},draw=SynBlue,thick}]
\node[caja] (a) {\textbf{1. Registrar y clasificar}\quad Proyecto y Funcional: requisito, motivo y solicitante; defecto, precisión o cambio contractual.};
\node[caja,below=5mm of a] (b) {\textbf{2. Analizar impacto}\quad Arquitectura, Datos, Seguridad, Calidad y SRE: alcance, plazo, capacidad, costo, riesgo y servicio.};
\node[caja,below=5mm of b] (c) {\textbf{3. Resolver antes de ejecutar}\quad Comité Ejecutivo y formalización escrita de ambas partes. Rechazo o aplazamiento: registrar decisión y comunicar.};
\node[caja,below=5mm of c] (d) {\textbf{4. Actualizar y comprobar}\quad Proyecto: versiones de alcance, EDT, calendario y T-12; responsables: implementar y probar; CLIENTE: recepción aplicable.};
\draw[enlace] (a)--(b);\draw[enlace] (b)--(c);\draw[enlace] (c)--node[right,font=\fontsize{9}{11}\selectfont]{Aprobado}(d);
\end{tikzpicture}
\caption{Control integrado de una modificación contractual}\label{fig:sd6-cambios}
\FuenteFigura{Elaboración de Synaptix a partir de las Bases Administrativas, art.~72, p.~38, y las Bases Técnicas Transversales, RT-19.03, p.~33.}
\end{figure}

La secuencia impide ejecutar una modificación mientras sólo existe una propuesta. El Comité Ejecutivo aprueba antes de ejecutar y ambas partes formalizan por escrito. Proyecto actualiza las líneas base y comunica la versión. Se comprueban límite acumulado, objeto, garantías y calendario del artículo 72, manteniendo los importes fuera de la oferta técnica. La Figura describe el procedimiento ofrecido, sin acreditar aprobaciones obtenidas.

La ejecución de un cambio sin aprobación previa del Comité Ejecutivo y acuerdo escrito de ambas partes será de cargo y riesgo exclusivo de Synaptix y no dará derecho a pago adicional. El procedimiento CC-6 de T-9 bloquea la asignación no autorizada, conserva la decisión y el análisis de impacto y registra cualquier incumplimiento con su contención y corrección; la ejecución no se convierte por sí sola en autorización. \parencite[art.~72, numerales 1--5, p.~38]{sd6-ba}

Las adquisiciones se vinculan a paquetes y fechas de necesidad de SD7 y al inventario T-11. Arquitectura define compatibilidad, Seguridad verifica licencia y tratamiento de datos y Calidad comprueba recepción. Se registran proveedor, producto, dependencia, plazo, aceptación, sustitución y riesgo de suministro. La justificación de hacer o comprar y la autorización de subcontratación se desarrollan en SD12; el método controla interfaces y no presume acuerdos firmados. \parencite[arts.~71--73, pp.~37--38]{sd6-ba}

\subsection{Seguimiento cuantitativo y respuesta a desviaciones}
El informe mensual contiene entregables, estado de hitos, desviaciones, riesgos, incidencias y compromisos siguientes. El tablero enlaza resultados y actas; el avance se mide con valor ganado. Se emplean $SPI=EV/PV$ y $CPI=EV/AC$: $PV$ es valor presupuestado del trabajo planificado, $EV$ el del trabajo realizado y $AC$ su costo real. Deben usar idéntica base monetaria y corte: sustituir $AC$ por horas sin ponderación no acredita CPI. Las fórmulas se sustentan en la ficha oficial de NASA; sus umbrales de contratación no se trasladan a esta oferta. \parencites[RT-19.06--10, p.~34]{sd6-bt}[p.~1]{sd6-nasa}

Los importes se administrarán en el expediente contractual y económico autorizado. SD6 fija el procedimiento sin publicar precios ni valores que permitan inferir la oferta, conforme al artículo 50.2. Los paquetes ganan valor mediante resultados verificables; T-9 establece hitos ponderados y tratamiento de servicios periódicos. Con denominador cero, el índice se declara no calculable y se explica la causa. Un índice agregado no reemplaza el control de ruta crítica.

Una desviación superior al 10\,\% en un hito se comunica dentro de cinco días hábiles de detectada, con causa, impacto y recuperación. Se comparan fecha proyectada, holgura y capacidad sin alterar retrospectivamente la línea base para ocultar retrasos. La gestión del riesgo aplica ISO 31000:2018 mediante el plan de SD8 y el registro T-16: contexto, identificación, análisis, evaluación, tratamiento, comunicación y seguimiento. En cada Comité de Proyecto se revisan responsable, respuesta, plazo, disparador y riesgo residual; un riesgo materializado se enlaza al incidente sin borrar su historial. \parencites[RT-19.04/10, pp.~33--34]{sd6-bt}{sd6-iso31000}

\subsection{Protección de la operación y evidencia de aplicación}
La marina impone congelamiento total del 15 de diciembre al 15 de marzo, protección del varadero en abril--mayo y septiembre--noviembre, y congelamientos locales en 14 fechas náuticas. Un aviso de marejadas suspende inmediatamente toda actividad programada del proyecto, incluida la construcción aislada o remota. Proyecto registra suspensión, trabajo afectado, capacidad y consumo de reserva; sólo reprograma tras cesar la condición y comprobar una ventana admisible. Las reservas de SD7 no se vuelven a sumar ni autorizan desplazar hitos contractuales. La atención de una emergencia real conserva su procedimiento de incidente y no habilita trabajo programado. Los mantenimientos programados requieren además aviso al CLIENTE con al menos diez días hábiles. \parencites[cap.~10, restricción 12, p.~26; numeral 13.2, p.~29; cap.~15, RT-10.05, p.~34]{sd6-c07}[art.~20, p.~15]{sd6-ba}[RT-10.05, p.~22]{sd6-bt}

El expediente enlaza versiones de línea base, capacidad, actas, riesgos, cambios, aceptaciones y reportes. T-9 define campos, responsables y cierre. La comprobación del método exigirá un cambio trazable desde solicitud hasta actualización documental, un corte de avance reproducible desde paquetes y la comprobación de que el trabajo quincenal cabe en la capacidad autorizada.
